\documentclass[11pt,a4paper]{article}
\usepackage[utf8]{inputenc}
\usepackage[T1]{fontenc}
\usepackage[spanish,es-nodecimaldot]{babel}
\usepackage{lmodern,microtype,amsmath,amssymb,graphicx,booktabs,array,placeins}
\usepackage[margin=2.4cm,headheight=14pt]{geometry}
\usepackage{xcolor,hyperref,fancyhdr}
\definecolor{ccdenavy}{HTML}{0A1E40}
\definecolor{ccdeblue}{HTML}{0077B6}
\hypersetup{colorlinks=true,linkcolor=ccdenavy,urlcolor=ccdeblue,citecolor=ccdeblue,pdftitle={Estructura productiva y recuperacion departamental en el Peru},pdfauthor={Circulo de Ciencia de Datos y Econometria (CCDE)}}
\graphicspath{{../figures/}}
% Generado por src/publicar.py; no editar a mano.
\newcommand{\BaseObs}{428}
\newcommand{\RatesObs}{384}
\newcommand{\LevelsObs}{168}
\newcommand{\Departments}{24}
\newcommand{\Sectors}{11}
\newcommand{\Zeros}{0}
\newcommand{\PreStart}{2015}
\newcommand{\PreEnd}{2019}
\newcommand{\RateStart}{2008}
\newcommand{\RateEnd}{2023}
\newcommand{\LevelStart}{2017}
\newcommand{\LevelEnd}{2023}
\newcommand{\BaseYear}{2019}
\newcommand{\ShockYear}{2020}
\newcommand{\PostStart}{2021}
\newcommand{\PostEnd}{2023}
\newcommand{\DeltaZero}{0.05}
\newcommand{\BaseStart}{2008}
\newcommand{\BaseEnd}{2025}
\newcommand{\GapReplicas}{30}
\newcommand{\StabilityReplicas}{100}
\newcommand{\SubsamplePct}{80}
\newcommand{\PcaOne}{51.65}
\newcommand{\PcaTwo}{20.02}
\newcommand{\PcaCumulative}{71.68}
\newcommand{\PcaThree}{87.39}
\newcommand{\SdRatio}{4.3}
\newcommand{\LeveneP}{0.0137}
\newcommand{\FlignerP}{0.0009}
\newcommand{\MacroMiningCount}{10}
\newcommand{\MacroOtherCount}{14}
\newcommand{\GapChoice}{1}
\newcommand{\SdEnclave}{21.3}
\newcommand{\SdAgro}{4.9}
\newcommand{\SpecCount}{48}
\newcommand{\ConfidencePct}{95}
\newcommand{\SilMacro}{0.498}
\newcommand{\SilTypology}{0.407}
\newcommand{\JaccardWard}{85.0}
\newcommand{\JaccardKmeans}{87.7}
\newcommand{\GapTypology}{0.231}
\newcommand{\WardSmallCount}{3}
\newcommand{\WardLargeCount}{21}
\newcommand{\ROne}{0.145}
\newcommand{\NOne}{384}
\newcommand{\RTwo}{0.141}
\newcommand{\NTwo}{384}
\newcommand{\RThree}{0.142}
\newcommand{\NThree}{384}
\newcommand{\RFour}{0.277}
\newcommand{\NFour}{383}
\newcommand{\RFive}{0.355}
\newcommand{\NFive}{168}
\newcommand{\RSix}{0.345}
\newcommand{\NSix}{168}
\newcommand{\ScarBeta}{-10.92}
\newcommand{\ScarP}{0.0029}
\newcommand{\ScarLow}{-18.05}
\newcommand{\ScarHigh}{-3.79}
\newcommand{\CovidEnclaveBeta}{-6.57}
\newcommand{\CovidEnclaveP}{0.0495}
\newcommand{\ContinuousBeta}{-2.76}
\newcommand{\ContinuousP}{0.0747}
\newcommand{\MacroCovidBeta}{-14.45}
\newcommand{\MacroCovidP}{$<0.0001$}
\newcommand{\AgroFinal}{104.9}
\newcommand{\AgroCount}{3}
\newcommand{\ServicesFinal}{103.1}
\newcommand{\ServicesCount}{11}
\newcommand{\MiningFinal}{106.6}
\newcommand{\MiningCount}{7}
\newcommand{\EnclaveFinal}{97.6}
\newcommand{\EnclaveCount}{3}
\newcommand{\EnclaveFirst}{95.9}
\newcommand{\EnclaveSecond}{94.6}
\newcommand{\EnclaveThird}{97.6}
\newcommand{\ExtremeGrowth}{141.2}

\pagestyle{fancy}
\fancyhf{}
\fancyhead[L]{\small\textcolor{ccdenavy}{CCDE | Documento de trabajo}}
\fancyhead[R]{\small Estructura productiva regional}
\fancyfoot[C]{\thepage}
\renewcommand{\headrulewidth}{0pt}
\setlength{\parskip}{.35em}
\setlength{\parindent}{1.2em}
\setlength{\emergencystretch}{2em}
\renewcommand{\arraystretch}{1.15}
\newcommand{\nota}[1]{\par\vspace{.2em}{\footnotesize\textit{Nota:} #1\par}}
\newcommand{\figura}[4]{\begin{figure}[htbp]\centering\includegraphics[width=\linewidth,height=.62\textheight,keepaspectratio]{#1}\caption{#2}\label{#3}\nota{#4}\end{figure}}
\title{\textcolor{ccdenavy}{Estructura productiva y recuperación departamental en el Perú}\\[.35em]\large Análisis composicional, volatilidad y persistencia del rezago tras el COVID-19}
\author{Círculo de Ciencia de Datos y Econometría (CCDE)}
\date{Documento de trabajo · Octubre de 2026}
\begin{document}
\maketitle
\begin{center}\includegraphics[width=1.75cm]{../assets/ccde-logo-positivo.png}\end{center}
\begin{abstract}
La recuperación posterior a 2020 fue heterogénea entre los departamentos del Perú. Estudiamos su relación con la estructura productiva previa mediante análisis composicional y modelos de panel con efectos fijos departamentales. Las participaciones de \Sectors\ actividades en \PreStart--\PreEnd\ se transforman mediante logaritmos de cocientes y se resumen con componentes principales. Una tipología Ward distingue cuatro perfiles: baja extracción, servicios y manufactura, minería diversificada y enclaves mineros. La validación muestra criterios en tensión: K=2 presenta mayor silueta, K=4 mayor estabilidad entre las particiones evaluadas y la regla Gap de un error estándar selecciona K=\GapChoice. En un panel del crecimiento del valor agregado bruto (VAB) real de \RateStart--\RateEnd\ (\RatesObs\ observaciones), la dispersión limita la precisión de algunos contrastes. En niveles indexados a \BaseYear\ = 100, con \LevelsObs\ observaciones de \LevelStart--\LevelEnd, los enclaves registran un cambio relativo de \ScarBeta\ puntos en \PostStart--\PostEnd\ frente al perfil diversificado/agro y al periodo \LevelStart--\PreEnd\ (IC \ConfidencePct\%: [\ScarLow, \ScarHigh]; $p=\ScarP$). Su índice medio cierra \LevelEnd\ en \EnclaveFinal. Los resultados documentan persistencia observada del rezago, sin identificar un efecto causal exclusivo de la pandemia ni demostrar una pérdida permanente de producción.
\end{abstract}
\noindent\textbf{Palabras clave:} economía regional, VAB departamental, datos composicionales, clustering, efectos fijos, COVID-19.\par
\noindent\textbf{JEL:} C23, C38, R11.
\clearpage
\section{Introducción}
Una tasa positiva de crecimiento posterior a una recesión no implica recuperar el nivel de actividad previo. Esta distinción es relevante al comparar economías regionales con estructuras productivas diferentes. Un departamento dependiente de un proyecto extractivo puede registrar variaciones pronunciadas de producción por cambios en capacidad o interrupciones operativas; otro con mayor participación de servicios puede mostrar tasas más estables. La comparación de tasas anuales ofrece por ello una imagen distinta de la comparación de niveles normalizados.

Este documento examina cómo se relaciona la estructura productiva previa a 2020 con la volatilidad del crecimiento y la recuperación observada hasta \LevelEnd. La pregunta empírica es si los perfiles productivos presentan cambios diferentes durante el choque de \ShockYear\ y el periodo \PostStart--\PostEnd. El análisis separa tres operaciones: construir una clasificación con información previa, describir la dispersión histórica y estimar cambios relativos dentro de cada departamento.

La contribución es integrar esas operaciones en un flujo reproducible. Tratamos las participaciones sectoriales como composiciones que suman 100\%, en lugar de interpretarlas como variables libres e independientes. Comparamos cuatro familias de agrupamiento y tres representaciones de los datos. Finalmente, contrastamos una misma estructura de interacciones en modelos cuya dependiente es el crecimiento anual o el nivel de VAB real indexado.

El principal resultado es que el perfil minero/enclave permanece por debajo de su nivel medio de 2019 durante 2021--2023, mientras los demás perfiles terminan por encima de ese umbral. El modelo en índices estima un cambio relativo negativo para ese grupo. Su interpretación requiere cautela: hay solamente \EnclaveCount\ departamentos enclave y la clasificación utiliza una reducción de dimensión que no conserva toda la geometría composicional. La evidencia describe una recuperación desigual; no permite atribuirla exclusivamente a la pandemia, a conflictos o a una ausencia de diversificación.

\section{Datos, unidad de medida y muestras}
\subsection{Fuente y armonización}
Utilizamos los 24 libros departamentales del INEI conservados en el repositorio. El cuadro 1 contiene el \textit{valor agregado bruto} a precios constantes de 2007 y el cuadro 2 su estructura porcentual por actividades. Las notas de los libros consultados señalan información disponible a junio de 2026. El portal del INEI ofrece el marco de las cuentas departamentales y publicaciones de referencia \cite{inei}; la versión exacta utilizada se identifica mediante los archivos locales y sus huellas SHA-256, sin suponer que una publicación enlazada reproduce la misma revisión.

La columna heredada \texttt{PIB\_total} de la base consolidada corresponde a la fila \textit{Valor Agregado Bruto}. La salida final utiliza \texttt{VAB\_total} y \texttt{VAB\_idx2019}. Esta distinción evita confundir VAB con PBI, que también incorpora impuestos netos sobre los productos. Las tasas y los índices calculados se refieren a la misma magnitud de VAB real.

La base original presenta 12 columnas sectoriales porque una actividad residual aparece bajo dos nombres según el departamento. Unificamos \textit{Otros Servicios} y \textit{Otros Actividades/1}; resultan \Sectors\ actividades comparables. La categoría residual incluye otros servicios y, según las notas de los libros, pesca y acuicultura. Por ello no debe interpretarse como una medición exclusiva del sector servicios. Verificamos el cierre de las participaciones a 100\% antes de calcular cocientes.

La unidad transversal es la consignada en cada libro departamental; se conserva la cobertura de Lima tal como figura en su archivo. No se construye una observación adicional para Callao ni se mezclan subunidades de Lima con el total del libro. Las medias por perfil otorgan el mismo peso a cada departamento: no representan el VAB agregado nacional ni un promedio ponderado por tamaño económico.

\subsection{Ventanas con propósitos diferentes}
La base local contiene \BaseObs\ observaciones entre \BaseStart\ y \BaseEnd. Los años 2008--2023 están completos para los \Departments\ departamentos; 2024 y 2025 cuentan con 22 unidades cada uno. Se conservan esos años en la base, pero se excluyen de las muestras de estimación para evitar que la cobertura incompleta altere las comparaciones.

La clasificación utiliza el promedio aritmético de participaciones de \PreStart--\PreEnd. Esta ventana precede al COVID-19 y evita clasificar departamentos con los propios resultados de la recuperación. No garantiza exogeneidad estricta: la composición puede relacionarse con tendencias previas, inversión o instituciones no observadas.

Los modelos de crecimiento utilizan \RateStart--\RateEnd, con \RatesObs\ observaciones balanceadas. La variante M4 retira una observación para sensibilidad y utiliza \NFour\ filas. Los modelos de niveles utilizan \LevelStart--\LevelEnd, con \LevelsObs\ observaciones y siete años por departamento. El año base es \BaseYear. Las tasas disponibles en la base consolidada están redondeadas a dos decimales; se verifican contra los niveles anuales disponibles.

\section{Metodología}
\subsection{Participaciones y geometría composicional}
Sea $x_i=(x_{i1},\ldots,x_{iD})$ la composición media del departamento $i$, con $D=\Sectors$, $x_{ij}\geq0$ y $\sum_jx_{ij}=100$. El cierre implica dependencia entre las partes: aumentar una participación requiere que otra disminuya. Trabajamos con logaritmos de cocientes, siguiendo el enfoque composicional de Aitchison \cite{aitchison}.

El algoritmo admite reemplazo multiplicativo de ceros \cite{martin}. Si $z_i$ es el número de partes nulas y $\delta=\DeltaZero$ puntos de participación, la transformación es
\begin{equation}
x^*_{ij}=\begin{cases}\delta,&x_{ij}=0,\\x_{ij}(100-z_i\delta)/100,&x_{ij}>0.\end{cases}
\end{equation}
Se aplica después de normalizar la composición a 100. En los promedios efectivamente utilizados se observan \textbf{\Zeros\ ceros}; por tanto, el reemplazo no modifica los datos. Valores pequeños de extracción en Lambayeque y San Martín no son ceros y no se sustituyen. La sensibilidad a $\delta$ es aquí una comprobación sin perturbación efectiva, no evidencia de robustez ante ceros estructurales.

Las coordenadas CLR son
\begin{equation}
\operatorname{clr}(x^*_{ij})=\log x^*_{ij}-\frac{1}{D}\sum_{h=1}^{D}\log x^*_{ih}.
\end{equation}
Sus filas suman cero y viven en un subespacio de dimensión $D-1$. Calculamos también coordenadas ILR con una base ortonormal de Helmert \cite{egozcue}. La distancia euclidiana entre vectores ILR coincide con la distancia de Aitchison entre composiciones.

Aplicamos PCA a CLR centrado, sin estandarizar columnas. Los primeros dos componentes explican \PcaCumulative\% de la varianza (\PcaOne\% y \PcaTwo\%); tres conservan \PcaThree\%. El primer eje se relaciona con extracción y el segundo contrasta manufactura con agropecuario, según sus cargas. El signo de cada componente es convencional. El mapa no superpone flechas de cargas en una escala arbitraria: las cargas se reportan por separado en el apéndice.

\subsection{Agrupamiento y validación}
Comparamos Ward, K-Means, PAM y mezclas gaussianas para K=2,3,4,5 sobre CLR+PCA2, CLR+PCA3 e ILR completo: \SpecCount\ configuraciones. K-Means utiliza 20 inicializaciones; PAM busca intercambios desde medoides iniciales; GMM utiliza cinco inicializaciones y se evalúa además mediante BIC. PAM es una búsqueda local, sin garantía de óptimo global. La semilla fija facilita reproducir los ejercicios.

Evaluamos silueta, Calinski--Harabasz y Davies--Bouldin. La silueta depende de la representación: un valor mayor tras reducir dimensión no demuestra que los componentes descartados sean ruido económico irrelevante. La partición principal de cuatro perfiles se utiliza por interpretabilidad y estabilidad relativa, manteniendo alternativas como sensibilidad.

Calculamos Gap con \GapReplicas\ referencias uniformes en la caja de las coordenadas PCA observadas \cite{gap}. Aplicamos la regla del primer K que cumple $\operatorname{Gap}(K)\geq\operatorname{Gap}(K+1)-s_{K+1}$. La estabilidad se aproxima con \StabilityReplicas\ submuestras del \SubsamplePct\% \textbf{sin reemplazo}: cada grupo original, restringido a las unidades seleccionadas, se compara con el grupo remuestreado de mayor Jaccard. Promediamos por grupo y luego por réplica. Este procedimiento es submuestreo, una alternativa discutida por Hennig \cite{hennig}; no es bootstrap ni un intervalo de confianza.

Los nombres K=4 se asignan por participación extractiva media ordenada. Son etiquetas descriptivas. En particular, \textit{diversificado/agro} incluye Lima, Lambayeque y San Martín; comparten baja extracción, sin que los tres sean economías predominantemente agrarias. Tampoco se deduce diversificación de una sola proporción minera; el HHI sectorial se exporta como descriptor complementario.

Distinguimos el corte Ward K=2 de una agregación económica de K=4. El primero separa \WardSmallCount\ y \WardLargeCount\ departamentos. La segunda agrupa los dos perfiles de mayor extracción en un bloque minero de \MacroMiningCount\ unidades y reúne los demás en otro de \MacroOtherCount. Se utiliza en M2 y no equivale al corte Ward de mayor silueta.

\subsection{Volatilidad, crecimiento e índice de nivel}
El crecimiento anual y el índice respecto a 2019 son
\begin{equation}
g_{it}=100\left(\frac{VAB_{it}}{VAB_{i,t-1}}-1\right),\qquad I_{it}=100\frac{VAB_{it}}{VAB_{i,2019}}.
\end{equation}
Una caída de 12\% seguida de un rebote de 9\% produce $100(0.88)(1.09)=95.92$. Es un ejemplo aritmético, no una trayectoria estimada. Las métricas informan dimensiones distintas.

La desviación estándar por perfil utiliza las observaciones departamento-año de $g_{it}$. Incluye heterogeneidad entre unidades, variación temporal, pandemia y rebote; no equivale a volatilidad idiosincrática tras controlar choques comunes. Reportamos Levene centrado en la mediana y Fligner--Killeen como diagnósticos descriptivos. Sus valores $p$ convencionales no ajustan la dependencia serial o transversal del panel.

\subsection{Modelos de efectos fijos}
Para los modelos categóricos estimamos
\begin{equation}\begin{split}
y_{it}={}&\alpha_i+\beta_C C_t+\beta_P P_t\\
&+\sum_{k\ne r}\gamma_k C_tG_{ik}+\sum_{k\ne r}\theta_k P_tG_{ik}+u_{it},
\end{split}\end{equation}
donde $C_t=1[t=2020]$, $P_t=1[2021\leq t\leq2023]$ y $r$ es diversificado/agro. La dependiente es $g_{it}$ en M1 y M4 e $I_{it}$ en M5. Los efectos fijos absorben los niveles constantes departamentales y los indicadores de grupo, que no varían en el tiempo.

M2 sustituye los cuatro perfiles por la agregación económica de dos grupos. M3 y M6 sustituyen las categorías por $z_i$, participación extractiva previa estandarizada con la media y desviación estándar muestral de los 24 departamentos. M3 usa crecimiento y M6 índice. M4 replica M1 excluyendo Apurímac 2016. La tabla \ref{tab:modelos} resume las especificaciones.

\begin{table}[htbp]\centering\small\begin{tabular}{llrrr}
\toprule
Modelo & Especificación & Periodo & Obs. & $R^2_w$ \\
\midrule
M1 & Tasas, cuatro perfiles & 2008--2023 & 384 & 0.145 \\
M2 & Tasas, dos macrogrupos & 2008--2023 & 384 & 0.141 \\
M3 & Tasas, minería continua & 2008--2023 & 384 & 0.142 \\
M4 & Tasas, sin Apurímac 2016 & 2008--2023 & 383 & 0.277 \\
M5 & Índice, cuatro perfiles & 2017--2023 & 168 & 0.355 \\
M6 & Índice, minería continua & 2017--2023 & 168 & 0.345 \\
\bottomrule
\end{tabular}

\caption{Especificaciones y muestras}\label{tab:modelos}
\nota{Todas incluyen efectos fijos departamentales. M4 retira una observación. El $R^2$ within no permite ordenar modelos con diferente dependiente y muestra.}\end{table}

Los coeficientes $\theta_k$ comparan el cambio entre \PostStart--\PostEnd\ y el periodo omitido con el cambio del grupo de referencia. En M5, el periodo omitido es \LevelStart--\PreEnd; en M1 comprende los años previos a 2020 de su ventana más larga. No es una brecha de nivel de un solo año ni la distancia frente a una trayectoria contrafactual sin pandemia. La base 2019 iguala el índice de todos los departamentos en ese año, pero no iguala sus trayectorias previas.

Los errores estándar se agrupan por departamento mediante \texttt{PanelOLS}; los intervalos se obtienen con \texttt{conf\_int} y la distribución t de la versión registrada del estimador \cite{linearmodels}. Se conserva el ajuste por regresores (\texttt{debiased=True}) y la configuración predeterminada sin ajuste adicional por número de grupos. Hay \Departments\ grupos de inferencia y solo \EnclaveCount\ departamentos enclave. No se incluyen efectos fijos anuales completos: los indicadores comunes recogen solo los regímenes definidos. El diseño no garantiza tendencias paralelas, ausencia de confusión temporal o asignación exógena a una estructura productiva.
\FloatBarrier

\section{Resultados}
\subsection{Tipología interpretable y criterios en tensión}
El mapa CLR+PCA2 distingue perfiles principalmente a lo largo del gradiente extractivo (figura \ref{fig:pca}). Las composiciones medias muestran diferencias entre los dos perfiles mineros en manufactura y otras actividades (figura \ref{fig:perfiles} y tabla \ref{tab:perfiles}). El grupo enclave incluye Apurímac, Huancavelica y Pasco; los demás miembros se enumeran en el apéndice.

\figura{fig1_pca_clusters.pdf}{Posiciones departamentales en CLR+PCA2}{fig:pca}{Ward K=4. Puntos y etiquetas usan colores por perfil; los ejes reportan varianza CLR explicada. Las cargas sectoriales se muestran en el apéndice. Elaboración CCDE con datos INEI.}
\figura{fig3_perfiles_sectoriales.pdf}{Composición sectorial media por perfil}{fig:perfiles}{Promedios simples entre departamentos de \PreStart--\PreEnd; cada barra suma 100\%. Otras actividades reúne las ramas no representadas individualmente e incluye pesca contenida en la categoría residual de origen.}
\begin{table}[htbp]\centering\small\begin{tabular}{lrrrr}
\toprule
Perfil & n & Extracción & Manufactura & Agro \\
\midrule
Diversificado/Agro & 3 & 0.9 & 12.7 & 12.3 \\
Servicios/Manufactura & 11 & 12.2 & 10.6 & 13.9 \\
Minero/Diversificado & 7 & 37.4 & 12.2 & 5.3 \\
Minero/Enclave & 3 & 41.9 & 1.8 & 8.1 \\
\bottomrule
\end{tabular}

\caption{Participaciones seleccionadas por perfil, en porcentaje}\label{tab:perfiles}
\nota{Las columnas seleccionadas no agotan la composición. Los nombres son etiquetas relativas derivadas de la participación extractiva.}\end{table}

Ward K=2 alcanza silueta \SilMacro, frente a \SilTypology\ en K=4. Ward y K-Means coinciden en K=4 sobre PCA2; su Jaccard medio es \JaccardWard\% y \JaccardKmeans\%, respectivamente. Sin embargo, el máximo local de Gap en K=4 (\GapTypology) no determina la selección formal: la regla de un error estándar elige K=\GapChoice. No corresponde describir K=4 como óptimo único o clasificación inequívoca (figura \ref{fig:validacion}).

PCA3 y la geometría ILR completa modifican la partición, con los ARI reportados en el apéndice. La sensibilidad limita las inferencias que toman el grupo estimado como una característica observada sin error. PCA2 conserva aproximadamente siete décimos de la variación CLR; el resto puede contener información económicamente relevante.

\figura{fig2_validacion_clustering.pdf}{Separación, Gap y estabilidad}{fig:validacion}{Paneles independientes para evitar doble eje. Silueta con origen cero, Jaccard en escala 0--1 y Gap con incertidumbre de Monte Carlo. La franja señala K=4; no es un intervalo de confianza.}
\FloatBarrier

\subsection{Dispersión y un episodio extremo}
La desviación estándar del crecimiento agrupado es \SdEnclave\ puntos porcentuales en enclave y \SdAgro\ en diversificado/agro: una razón de \SdRatio\ entre \textbf{desviaciones estándar}. No es una razón de varianzas, la cual sería el cuadrado de esa cifra. Levene arroja $p=\LeveneP$ y Fligner--Killeen $p=\FlignerP$. La dependencia del panel impide tratar esos valores como prueba definitiva de un régimen estructural de volatilidad.

Apurímac registra un aumento de \ExtremeGrowth\% en 2016. Es compatible con un gran cambio de capacidad extractiva y se conserva como observación económica real. Sin información adicional de proyectos no se atribuye exclusivamente a una operación concreta. La figura \ref{fig:volatilidad} muestra todas las observaciones en un eje lineal continuo y compara la dispersión antes y después de excluir ese punto como sensibilidad.

El $R^2$ within de M1 pasa de \ROne\ a \RFour\ en M4. Esto muestra la influencia de la observación en un modelo con pocos regresores de tiempo; no demuestra que retirarla mejore la especificación económica. M4 usa otra muestra y se interpreta como sensibilidad, manteniendo M1 como referencia del panel completo.

\figura{fig4_volatilidad_estructural.pdf}{Crecimiento y sensibilidad de la dispersión}{fig:volatilidad}{Panel \RateStart--\RateEnd. Puntos departamento-año; cajas con mediana y rango intercuartílico. Las barras parten de cero. La exclusión de M4 afecta una sola celda y no modifica la fuente.}
\FloatBarrier

\subsection{Contrastes en tasas de crecimiento}
M1--M4 se presentan en la tabla \ref{tab:tasas}. Los indicadores de choque recogen cambios del crecimiento respecto al periodo previo; no son directamente la tasa de 2020 de todos los departamentos. Los efectos principales dependen de la categoría de referencia o de la media del predictor estandarizado.

Los contrastes COVID por perfil de M1 son imprecisos, especialmente en los grupos extractivos. Un intervalo que contiene cero indica que la muestra no distingue con precisión algunos cambios relativos; no establece igualdad de impactos ni demuestra un error de tipo II. Tampoco identifica como causa de la imprecisión la simultaneidad del confinamiento. Dispersión, tamaños de grupo y especificación son explicaciones posibles que requieren nuevos contrastes.

\begin{table}[htbp]\centering\footnotesize\setlength{\tabcolsep}{4pt}
\begin{tabular}{lrrrr}
\toprule
Regresor & M1 & M2 & M3 & M4 \\
\midrule
Constante & \shortstack{4.58 \\ (0.22)} & \shortstack{4.58 \\ (0.23)} & \shortstack{4.58 \\ (0.23)} & \shortstack{4.10 \\ (0.17)} \\
Choque 2020 & \shortstack{-12.41 \\ (1.90)} & \shortstack{-14.45 \\ (0.89)} & \shortstack{-14.70 \\ (1.17)} & \shortstack{-12.41 \\ (1.90)} \\
Periodo 2021--2023 & \shortstack{-0.61 \\ (0.64)} & \shortstack{0.26 \\ (0.60)} & \shortstack{0.35 \\ (1.01)} & \shortstack{-0.61 \\ (0.64)} \\
2020 $\times$ Servicios/Manufactura & \shortstack{-2.60 \\ (2.12)} & --- & --- & \shortstack{-2.60 \\ (2.12)} \\
2020 $\times$ Minero/Diversificado & \shortstack{-1.22 \\ (3.78)} & --- & --- & \shortstack{-1.22 \\ (3.78)} \\
2020 $\times$ Minero/Enclave & \shortstack{-5.93 \\ (4.32)} & --- & --- & \shortstack{-2.12 \\ (2.79)} \\
2021--2023 $\times$ Servicios/Manufactura & \shortstack{1.11 \\ (0.97)} & --- & --- & \shortstack{1.11 \\ (0.97)} \\
2021--2023 $\times$ Minero/Diversificado & \shortstack{2.46 \\ (2.18)} & --- & --- & \shortstack{2.46 \\ (2.18)} \\
2021--2023 $\times$ Minero/Enclave & \shortstack{-2.11 \\ (5.43)} & --- & --- & \shortstack{1.70 \\ (2.78)} \\
2020 $\times$ Macro minero & --- & \shortstack{-0.59 \\ (2.79)} & --- & --- \\
2021--2023 $\times$ Macro minero & --- & \shortstack{0.22 \\ (2.34)} & --- & --- \\
2020 $\times$ Minería estandarizada & --- & --- & \shortstack{-1.77 \\ (1.12)} & --- \\
2021--2023 $\times$ Minería estandarizada & --- & --- & \shortstack{-0.15 \\ (1.40)} & --- \\
Observaciones & 384 & 384 & 384 & 383 \\
$R^2$ within & 0.145 & 0.141 & 0.142 & 0.277 \\
\bottomrule
\end{tabular}

\caption{Crecimiento anual del VAB real, M1--M4}\label{tab:tasas}
\nota{Coeficientes en puntos porcentuales; errores estándar agrupados por departamento entre paréntesis. M1 y M4: referencia diversificado/agro; M2: macrogrupo diversificado; M3: minería estandarizada. Valores $p$ e intervalos completos en la tabla procesada.}\end{table}
\FloatBarrier

\subsection{Recuperación y alcance del rezago}
La figura \ref{fig:trayectoria} muestra promedios simples de índices departamentales. La senda enclave es \EnclaveFirst, \EnclaveSecond\ y \EnclaveThird\ entre 2021 y 2023, inferior a 100 en los tres años. Al cierre, los otros perfiles registran \AgroFinal\ (diversificado/agro), \ServicesFinal\ (servicios/manufactura) y \MiningFinal\ (minero diversificado). Son diferencias descriptivas; la media de índices no equivale al índice de la suma del VAB de cada grupo.

M5 estima para enclave un cambio relativo de \CovidEnclaveBeta\ puntos en 2020 ($p=\CovidEnclaveP$) y de \ScarBeta\ puntos en \PostStart--\PostEnd\ ($p=\ScarP$), respecto al perfil de referencia y a 2017--2019 (tabla \ref{tab:niveles}). El segundo intervalo al \ConfidencePct\% es [\ScarLow, \ScarHigh]. La normalización fuerza igualdad en 2019; diferencias en las sendas de 2017 y 2018 también inciden en el contraste entre periodos. Por ello no se interpreta como una pérdida absoluta de \ScarBeta\ puntos desde 2019.

La interacción continua de recuperación de M6 es \ContinuousBeta\ puntos por una desviación estándar adicional de extracción previa ($p=\ContinuousP$). Su dirección coincide con el contraste categórico, pero el intervalo convencional al 95\% incluye cero. No confirma por sí misma una relación causal o un gradiente preciso. La figura \ref{fig:forest} compara M1, M4 y M5 con escalas independientes para tasas e índices y unidades explícitas.

\figura{fig5_trayectoria_vab_indexado.pdf}{VAB real indexado, \BaseYear\ = 100}{fig:trayectoria}{Promedios sin ponderar de índices departamentales de \LevelStart--\LevelEnd. La línea horizontal señala el nivel de 2019 y la banda vertical 2020. Se omiten bandas de error estándar selectivas que no representan inferencia para todos los grupos.}
\begin{table}[htbp]\centering\small\begin{tabular}{lrr}
\toprule
Regresor & M5 & M6 \\
\midrule
Constante & \shortstack{97.71 \\ (0.94)} & \shortstack{97.71 \\ (0.93)} \\
Choque 2020 & \shortstack{-5.14 \\ (1.60)} & \shortstack{-7.83 \\ (1.44)} \\
Periodo 2021--2023 & \shortstack{7.17 \\ (1.63)} & \shortstack{4.19 \\ (1.81)} \\
2020 $\times$ Servicios/Manufactura & \shortstack{-1.77 \\ (1.82)} & --- \\
2020 $\times$ Minero/Diversificado & \shortstack{-3.63 \\ (4.96)} & --- \\
2020 $\times$ Minero/Enclave & \shortstack{-6.57 \\ (3.32)} & --- \\
2021--2023 $\times$ Servicios/Manufactura & \shortstack{-1.45 \\ (2.25)} & --- \\
2021--2023 $\times$ Minero/Diversificado & \shortstack{-3.26 \\ (5.71)} & --- \\
2021--2023 $\times$ Minero/Enclave & \shortstack{-10.92 \\ (3.61)} & --- \\
2020 $\times$ Minería estandarizada & --- & \shortstack{-2.39 \\ (1.40)} \\
2021--2023 $\times$ Minería estandarizada & --- & \shortstack{-2.76 \\ (1.54)} \\
Observaciones & 168 & 168 \\
$R^2$ within & 0.355 & 0.345 \\
\bottomrule
\end{tabular}

\caption{Índice de VAB real, M5--M6}\label{tab:niveles}
\nota{Errores estándar agrupados por departamento entre paréntesis. M5: referencia diversificado/agro; M6: extracción estandarizada. Periodo temporal omitido: \LevelStart--\PreEnd.}\end{table}
\figura{fig6_forest_plot_modelos.pdf}{Cambios relativos en tasas e índices}{fig:forest}{IC al \ConfidencePct\% obtenidos del estimador con errores agrupados; se muestran sin truncar. Referencia: diversificado/agro. Las ventanas temporales de tasas y niveles difieren.}
\FloatBarrier

\section{Discusión y limitaciones}
Denominamos \textit{cicatriz} al rezago observado del nivel de actividad durante el horizonte disponible. No implica demostrar histéresis permanente, pérdida de producto potencial ni distancia frente a una senda sin pandemia. Un índice inferior a 100 indica producción menor que en el año base; uno superior a 100 tampoco descarta pérdida frente a un contrafactual creciente.

Hay cuatro límites principales. Primero, la clasificación se estima sobre 24 composiciones y depende del número de componentes y grupos. Revisar varias métricas antes de elegir una partición introduce discrecionalidad. Los cambios con PCA3 e ILR completo muestran que esa incertidumbre importa. Los errores de las regresiones no propagan la incertidumbre de clasificación ni posibles errores de medición sectorial.

Segundo, los efectos fijos controlan factores constantes, pero no tendencias específicas, choques anuales no capturados o cambios simultáneos de proyectos e instituciones. Clasificar con información previa evita utilizar los resultados de recuperación al asignar grupos, pero no garantiza exogeneidad. La igualdad mecánica en 2019 tampoco verifica tendencias paralelas. Un estudio de eventos con efectos anuales y una estrategia de identificación explícita son extensiones, no validaciones ya realizadas.

Tercero, el perfil enclave comprende tres unidades. Los \Departments\ grupos de inferencia son pocos para confiar irrestrictamente en aproximaciones asintóticas; la precisión puede depender de departamentos concretos. Correcciones apropiadas para pocos grupos, un bootstrap de conglomerados compatible con el diseño y exclusiones por departamento serían verificaciones adicionales. El submuestreo de clustering no satisface esas necesidades econométricas.

Cuarto, la comparación es territorial y sectorial, no distributiva. No permite concluir sobre empleo, bienestar, ingresos fiscales, encadenamientos o eficacia del canon. Profundizar vínculos productivos o gestionar riesgos requiere datos y evaluaciones específicas. La aplicación inmediata es un tablero que distinga crecimiento, niveles recuperados y exposición sectorial, sin convertir asociaciones en recomendaciones causales.

\section{Conclusiones}
La estructura previa permite organizar los departamentos en perfiles útiles para comparar desempeño. K=4 describe diferencias, pero los criterios no convergen en un único número de grupos y la clasificación cambia con la representación. El corte Ward K=2 y la agregación económica minero/diversificado deben mantenerse separados.

Combinar tasas e índices muestra por qué el rebote no resume la recuperación. Los enclaves presentan mayor dispersión y su índice medio permanece bajo el nivel de 2019 hasta 2023. El contraste de efectos fijos es consistente con un cambio relativo desfavorable en 2021--2023, sujeto a los límites de muestra e identificación descritos.

El siguiente paso es contrastar mecanismos con inversión, operaciones y conflictividad, ampliar el horizonte con cobertura comparable y evaluar diseños alternativos. Hasta entonces, el resultado es evidencia de heterogeneidad y persistencia observada del rezago departamental.

\section{Reproducibilidad}
El repositorio conserva los libros fuente y la base consolidada. Los supuestos se centralizan en \texttt{config/analisis.toml}. El comando \texttt{python src/generar\_analisis\_completo.py} regenera composiciones, coordenadas, métricas, etiquetas, panel, estimaciones, figuras y tablas LaTeX. Este documento y la presentación importan las mismas cifras generadas, sin transcribir coeficientes manualmente.

El manifiesto \texttt{data/processed/resultados.json} registra cobertura, supuestos, versiones y huellas SHA-256 de entrada, fuentes y salidas procesadas. Las versiones utilizadas se fijan en \texttt{requirements.txt}. Los cuadernos cargan las salidas finales y no duplican modelos. El comando \texttt{python src/compilar.py} compila artículo y presentación, y actualiza PowerPoint desde las páginas finales del PDF. Las verificaciones numéricas y estructurales se ejecutan con \texttt{python src/verificar.py}; rechazan cambios de configuración o archivos posteriores al cálculo.

\noindent Repositorio: \url{https://github.com/CCDE-DataLab/clusters-departamentales}.
\FloatBarrier

\begin{thebibliography}{9}
\bibitem{aitchison} Aitchison, J. (1986). \textit{The Statistical Analysis of Compositional Data}. Chapman \& Hall.
\bibitem{martin} Martín-Fernández, J. A., Barceló-Vidal, C. y Pawlowsky-Glahn, V. (2003). Dealing with zeros and missing values in compositional data sets using nonparametric imputation. \textit{Mathematical Geology}, 35, 253--278. \url{https://doi.org/10.1023/A:1023866030544}.
\bibitem{egozcue} Egozcue, J. J., Pawlowsky-Glahn, V., Mateu-Figueras, G. y Barceló-Vidal, C. (2003). Isometric logratio transformations for compositional data analysis. \textit{Mathematical Geology}, 35, 279--300. \url{https://doi.org/10.1023/A:1023818214614}.
\bibitem{gap} Tibshirani, R., Walther, G. y Hastie, T. (2001). Estimating the number of clusters in a data set via the gap statistic. \textit{Journal of the Royal Statistical Society, Series B}, 63(2), 411--423. \url{https://doi.org/10.1111/1467-9868.00293}.
\bibitem{hennig} Hennig, C. (2007). Cluster-wise assessment of cluster stability. \textit{Computational Statistics \& Data Analysis}, 52(1), 258--271. \url{https://doi.org/10.1016/j.csda.2006.11.025}.
\bibitem{inei} INEI. \textit{Producto Bruto Interno por Departamentos}. Portal de estadísticas y cuentas departamentales, base 2007. \url{https://m.inei.gob.pe/estadisticas/indice-tematico/producto-bruto-interno-por-departamentos-9089/}. Consulta: 28 de octubre de 2026. Libros utilizados en \texttt{data/raw}; véase el manifiesto de fuentes.
\bibitem{linearmodels} Sheppard, K. y colaboradores. \textit{linearmodels: Panel models. Documentación de PanelOLS}. \url{https://bashtage.github.io/linearmodels/panel/model.html}. Versión de ejecución registrada en el manifiesto.
\end{thebibliography}
\clearpage
\appendix
\section{Clasificación y sensibilidades}
\begin{table}[htbp]\centering\small\begin{tabular}{p{.27\linewidth}p{.65\linewidth}}
\toprule
Perfil & Departamentos \\
\midrule
Diversificado/Agro & Lambayeque, Lima, San Martín \\
Servicios/Manufactura & Amazonas, Ayacucho, Cajamarca, Huánuco, Ica, La Libertad, Loreto, Piura, Puno, Tumbes, Ucayali \\
Minero/Diversificado & Arequipa, Cusco, Junín, Madre de Dios, Moquegua, Tacna, Áncash \\
Minero/Enclave & Apurímac, Huancavelica, Pasco \\
\bottomrule
\end{tabular}

\caption{Miembros de los perfiles Ward K=4}\end{table}
\begin{table}[htbp]\centering\small\begin{tabular}{rrrrr}
\toprule
K & Silueta Ward & Gap & EE Gap & Jaccard Ward \\
\midrule
1 & --- & 0.356 & 0.145 & --- \\
2 & 0.498 & 0.201 & 0.120 & 0.625 \\
3 & 0.350 & 0.140 & 0.148 & 0.736 \\
4 & 0.407 & 0.231 & 0.123 & 0.850 \\
5 & 0.367 & 0.057 & 0.164 & 0.786 \\
\bottomrule
\end{tabular}

\caption{Validación de Ward sobre PCA2}
\nota{Jaccard de submuestreo; no es una probabilidad de pertenencia. La regla Gap de un error estándar considera K=1.}\end{table}
\begin{table}[htbp]\centering\small\begin{tabular}{lr}
\toprule
Representación / reemplazo & ARI frente a PCA2 \\
\midrule
CLR+PCA2 & 1.000 \\
CLR+PCA3 & 0.567 \\
ILR-full & 0.867 \\
delta=0.01 & 1.000 \\
delta=0.05 & 1.000 \\
delta=0.1 & 1.000 \\
\bottomrule
\end{tabular}

\caption{Sensibilidad de K=4}
\nota{ARI frente a CLR+PCA2. Cambiar delta no perturba la composición porque los promedios observados no contienen ceros.}\end{table}
\FloatBarrier
\figura{fig3b_jerarquia_ward.pdf}{Jerarquía Ward y cortes derivados de las fusiones}{fig:jerarquia}{Colores de nombres según K=4. K=2 separa \WardSmallCount\ y \WardLargeCount\ unidades; no es la agregación económica \MacroMiningCount/\MacroOtherCount. Los cortes se calculan entre alturas consecutivas.}
\FloatBarrier
\section{Cargas de componentes principales}
\begin{table}[htbp]\centering\small\begin{tabular}{p{.58\linewidth}rrr}
\toprule
Actividad & PC1 & PC2 & PC3 \\
\midrule
Agricultura, Ganadería, Caza y Silvicultura & -0.130 & -0.542 & -0.373 \\
Extracción de Petróleo, Gas y Minerales & 0.906 & 0.137 & -0.241 \\
Manufactura & -0.152 & 0.630 & 0.212 \\
Electricidad, Gas y Agua & 0.145 & -0.289 & 0.825 \\
Construcción & 0.026 & -0.175 & 0.039 \\
Comercio & -0.181 & 0.016 & -0.216 \\
Transporte, Almacen., Correo y Mensajería & -0.110 & 0.152 & -0.073 \\
Alojamiento y Restaurantes & -0.090 & 0.169 & -0.023 \\
Telecom. y Otros Serv. de Información & -0.202 & 0.217 & -0.149 \\
Administración Pública y Defensa & -0.078 & -0.277 & 0.029 \\
Otros Actividades/1 & -0.133 & -0.038 & -0.029 \\
\bottomrule
\end{tabular}

\caption{Cargas PCA de las actividades armonizadas}
\nota{CLR centrado, sin estandarizar columnas. Los signos pueden invertirse sin alterar la geometría. Las cargas no son coeficientes de una regresión del VAB.}\end{table}

\end{document}
